\documentclass[UTF8,a4paper,12pt]{ctexart}

\usepackage{amsmath,amssymb,bm}
\usepackage{booktabs}
\usepackage{geometry}
\usepackage{array}
\usepackage{float}
\usepackage{graphicx}
\usepackage{caption}
\usepackage{hyperref}

\geometry{left=2.4cm,right=2.4cm,top=2.4cm,bottom=2.4cm}
\setlength{\parindent}{2em}
\linespread{1.18}
\captionsetup{font=small}
\hypersetup{colorlinks=true,linkcolor=black,citecolor=black,urlcolor=blue}

\begin{document}

\section{问题一：预热阶段的径向热质传递模型}

\subsection{模型思路与基本假设}

热风干燥过程中，外界空气通过药材表面对流换热，热量进一步向内部传导；与此同时，
表面失水形成水分梯度并驱动内部水分向表面迁移。
中药材干燥研究通常将其概括为内部“导湿”和表面“给湿”的同步过程
\cite{wang2024}。
由于题目要求给出不同径向位置处的瞬时温度和水分浓度，必须保留空间分布，
因此采用一维圆柱坐标下的非稳态传热--传质模型。
Da Silva 等针对圆柱形物料建立了一维热质传递模型，
采用第三类边界和全隐式有限体积法求解\cite{dasilva2014}，
其几何结构和边界形式与本题具有较强对应性。

药材长度 $L=0.25~\mathrm m$、半径 $R=0.02~\mathrm m$，故
\[
\frac{L}{R}=12.5,\qquad
\frac{\text{两端总面积}}{\text{侧面积}}=\frac{R}{L}=0.08.
\]
因此以圆柱中段为代表区域，忽略端面附近的轴向梯度，仅考虑径向变化。
同时，
\[
Bi_T=\frac{h_TR}{k}\approx1.389,
\qquad
Bi_C=\frac{h_mR}{D(C_0)}\approx3.24,
\]
均明显大于集中参数模型常用的 $0.1$ 判据，
说明内部温度和水分梯度不可忽略。

问题一作如下简化：
\begin{enumerate}
    \item 药材视为均质、各向同性连续介质，初始温度和干基含水率均匀；
    \item 前 $1800~\mathrm s$ 内圆柱几何尺寸不变；
    \item $\rho,c_p,k$ 采用附录2给定常数；
    \item 附件1中的烘房水分浓度作为等效外界传质驱动力；
    \item 暂不考虑蒸发潜热、辐射换热及内部热源对温度场的反馈。
\end{enumerate}

其中第5条属于基础模型简化，并不意味着潜热在真实干燥过程中必然可以忽略；
已有研究表明，潜热对部分物料的温度预测可能产生明显影响\cite{dasilva2014}。

\subsection{主要参数}

\begin{table}[H]
\centering
\caption{问题一主要几何参数与材料参量}
\begin{tabular}{clcl}
\toprule
符号 & 物理意义 & 数值或关系式 & 单位\\
\midrule
$L$ & 药材长度 & $0.25$ & m\\
$R$ & 药材半径 & $0.02$ & m\\
$T_0$ & 初始温度 & $28$ & $^\circ$C\\
$C_0$ & 初始干基含水率 & $2.55$ & kg/kg\\
$\rho$ & 密度 & $820$ & kg/m$^3$\\
$c_p$ & 比热容 & $2600$ & J/(kg$\cdot$K)\\
$k$ & 导热系数 & $0.36$ & W/(m$\cdot$K)\\
$h_T$ & 对流换热系数 & $25$ & W/(m$^2\cdot$K)\\
$h_m$ & 等效对流传质系数 & $8\times10^{-7}$ & m/s\\
$D(C)$ & 有效水分扩散系数 & $7\times10^{-9}e^{-0.89/C}$ & m$^2$/s\\
\bottomrule
\end{tabular}
\end{table}

\subsection{控制方程与初边界条件}

\subsubsection{温度场}

根据能量守恒和 Fourier 导热定律，
\begin{equation}
\boxed{
\rho c_p\frac{\partial T}{\partial t}
=
\frac1r\frac{\partial}{\partial r}
\left(
rk\frac{\partial T}{\partial r}
\right)
},
\qquad 0<r<R.
\label{eq:heat}
\end{equation}

热扩散率为
\[
\alpha=\frac{k}{\rho c_p}
\approx1.6886\times10^{-7}\ \mathrm{m^2/s}.
\]

\subsubsection{水分场}

以干基含水率 $C=m_w/m_d$ 为状态变量，
采用有效扩散模型描述内部水分迁移：
\begin{equation}
\boxed{
\frac{\partial C}{\partial t}
=
\frac1r\frac{\partial}{\partial r}
\left[
rD(C)\frac{\partial C}{\partial r}
\right]
},
\qquad 0<r<R,
\label{eq:mass}
\end{equation}
其中
\[
D(C)=7\times10^{-9}
\exp\left(-\frac{0.89}{C}\right)
\ \mathrm{m^2/s}.
\]
由于 $D$ 随局部含水率变化，计算中保留其通量散度形式。

初始条件为
\[
T(r,0)=28^\circ\mathrm C,\qquad
C(r,0)=2.55.
\]

轴线处满足对称条件：
\begin{equation}
\left.\frac{\partial T}{\partial r}\right|_{r=0}=0,
\qquad
\left.\frac{\partial C}{\partial r}\right|_{r=0}=0.
\end{equation}

侧表面采用第三类边界：
\begin{equation}
-k\left.\frac{\partial T}{\partial r}\right|_{r=R}
=
h_T\left(T_s-T_\infty\right),
\end{equation}
\begin{equation}
-D(C_s)\left.\frac{\partial C}{\partial r}\right|_{r=R}
=
h_m\left(C_s-C_\infty\right).
\end{equation}

其中 $T_s=T(R,t)$、$C_s=C(R,t)$。
附件1中的烘房水分浓度与药材干基含水率并非严格相同的热力学浓度定义，
因此本文将其作为题设传质系数下的等效外界浓度驱动力。

附件1以 $60~\mathrm s$ 为间隔给出 $T_\infty$ 和 $C_\infty$，
计算中采用分段线性插值：
\begin{equation}
B(t)=B(\tau_j)+
\frac{t-\tau_j}{\tau_{j+1}-\tau_j}
\left[B(\tau_{j+1})-B(\tau_j)\right],
\end{equation}
其中 $B$ 表示 $T_\infty$ 或 $C_\infty$。

\subsection{数值求解方法}

由于模型包含圆柱几何项、时变 Robin 边界和非线性扩散系数，
采用径向有限体积法，并使用后向 Euler 全隐式时间推进。
第 $i$ 个控制体的体积和界面面积分别为
\begin{align}
V_i&=\pi L\left(r_{i+1/2}^2-r_{i-1/2}^2\right),\\
A_{i\pm1/2}&=2\pi Lr_{i\pm1/2}.
\end{align}

\IfFileExists{问题1_有限体积离散示意图_论文风格.png}{
\begin{figure}[H]
\centering
\includegraphics[width=0.82\textwidth]{问题1_有限体积离散示意图_论文风格.png}
\caption{圆柱药材径向有限体积离散示意图}
\end{figure}
}{
\begin{figure}[H]
\centering
\fbox{\parbox[c][3.8cm][c]{0.82\textwidth}{\centering
此处插入“问题1\_有限体积离散示意图\_论文风格.png”}}
\caption{圆柱药材径向有限体积离散示意图}
\end{figure}
}

离散后的温度方程为
\begin{equation}
\rho c_pV_i
\frac{T_i^{n+1}-T_i^n}{\Delta t}
=
A_{i-1/2}q_{i-1/2}^{n+1}
-
A_{i+1/2}q_{i+1/2}^{n+1},
\end{equation}
水分方程为
\begin{equation}
V_i
\frac{C_i^{n+1}-C_i^n}{\Delta t}
=
A_{i-1/2}f_{i-1/2}^{n+1}
-
A_{i+1/2}f_{i+1/2}^{n+1}.
\end{equation}

其中
\[
q=-k\frac{\partial T}{\partial r},
\qquad
f=-D(C)\frac{\partial C}{\partial r}.
\]
内部水分界面的扩散系数采用调和平均：
\begin{equation}
D_{i+1/2}
=
\frac{2D(C_i)D(C_{i+1})}
{D(C_i)+D(C_{i+1})}.
\end{equation}

对于最外层控制体，将半单元传递阻力与外部对流阻力串联：
\begin{equation}
q_R=
\frac{T_P-T_\infty}
{\dfrac{\Delta r}{2k}+\dfrac1{h_T}},
\qquad
T_s=T_\infty+\frac{q_R}{h_T},
\end{equation}
\begin{equation}
f_R=
\frac{C_P-C_\infty}
{\dfrac{\Delta r}{2D_s}+\dfrac1{h_m}},
\qquad
C_s=C_\infty+\frac{f_R}{h_m}.
\end{equation}

温度方程为线性三对角系统；
水分方程因 $D=D(C)$ 为非线性，在每个新时间层采用 Picard 迭代，
并以
\[
\max\left(
\|C^{(m+1)}-C^{(m)}\|_\infty,
|C_s^{(m+1)}-C_s^{(m)}|
\right)<10^{-10}
\]
作为收敛条件。最终计算中，每个时间层最多4次迭代即可收敛。

\subsection{数值可靠性检验}

最终采用
\[
N=1600,\qquad
\Delta r=1.25\times10^{-5}\ \mathrm m,
\qquad
\Delta t=0.015625\ \mathrm s.
\]

以题目规定的全部输出点为比较集合，
分别检验空间网格、时间步长和 Picard 容差的影响。

\begin{table}[H]
\centering
\caption{问题一数值可靠性检验}
\begin{tabular}{lccc}
\toprule
检验项目 & 对比设置 & $E_T$/($^\circ$C) & $E_C$/(kg/kg)\\
\midrule
空间网格细化
& $N=800\rightarrow1600$
& $1.02\times10^{-6}$
& $1.522\times10^{-5}$\\
时间步细化
& $0.03125\rightarrow0.015625$ s
& $2.986\times10^{-5}$
& $6.39\times10^{-6}$\\
Picard 容差收紧
& $10^{-8}\rightarrow10^{-10}$
& $0$
& $2.15\times10^{-13}$\\
\bottomrule
\end{tabular}
\end{table}

同时，最大相对离散平衡残差为
\[
\varepsilon_T^{\mathrm{bal}}
=3.919\times10^{-7},
\qquad
\varepsilon_C^{\mathrm{bal}}
=3.103\times10^{-11}.
\]
上述结果表明离散误差和非线性迭代误差均已足够小。
这里的平衡检验用于验证离散模型的数值闭合性，
并不替代真实内部温度和含水率观测所提供的实验验证。

\subsection{计算结果与物理解释}

\begin{table}[H]
\centering
\caption{30 min 内药材径向温度分布（$^\circ$C）}
\begin{tabular}{rrrrrr}
\toprule
时间/s & 0 cm & 0.5 cm & 1.0 cm & 1.5 cm & 2.0 cm\\
\midrule
100  & 28.0001 & 28.0003 & 28.0040 & 28.0327 & 28.1801\\
300  & 28.0408 & 28.0635 & 28.1514 & 28.3681 & 28.8489\\
600  & 28.4534 & 28.5361 & 28.8040 & 29.3159 & 30.1652\\
900  & 29.3243 & 29.4583 & 29.8755 & 30.6162 & 31.7304\\
1200 & 30.5427 & 30.7098 & 31.2224 & 32.1126 & 33.4277\\
1500 & 31.9957 & 32.1868 & 32.7660 & 33.7463 & 35.1203\\
1800 & 33.5754 & 33.7720 & 34.3642 & 35.3621 & 36.7856\\
\bottomrule
\end{tabular}
\end{table}

\begin{table}[H]
\centering
\caption{30 min 内药材径向水分浓度分布（kg/kg，干基）}
\begin{tabular}{rrrrrr}
\toprule
时间/s & 0 cm & 0.5 cm & 1.0 cm & 1.5 cm & 2.0 cm\\
\midrule
100  & 2.5500 & 2.5500 & 2.5500 & 2.5500 & 2.2470\\
300  & 2.5500 & 2.5500 & 2.5500 & 2.5492 & 2.0508\\
600  & 2.5500 & 2.5500 & 2.5500 & 2.5353 & 1.8770\\
900  & 2.5500 & 2.5500 & 2.5497 & 2.5045 & 1.7547\\
1200 & 2.5500 & 2.5500 & 2.5482 & 2.4646 & 1.6586\\
1500 & 2.5500 & 2.5499 & 2.5445 & 2.4206 & 1.5787\\
1800 & 2.5500 & 2.5497 & 2.5383 & 2.3755 & 1.5102\\
\bottomrule
\end{tabular}
\end{table}

\IfFileExists{问题1_径向温度分布_重算.png}{
\begin{figure}[H]
\centering
\includegraphics[width=0.82\textwidth]{问题1_径向温度分布_重算.png}
\caption{不同预热时刻的径向温度分布}
\end{figure}
}{}

\IfFileExists{问题1_径向水分浓度分布_重算.png}{
\begin{figure}[H]
\centering
\includegraphics[width=0.82\textwidth]{问题1_径向水分浓度分布_重算.png}
\caption{不同预热时刻的径向水分浓度分布}
\end{figure}
}{}

结果显示，温度变化已经深入圆柱主体区域，而水分变化仍主要集中在近表层。
这一差异可由特征尺度解释。

在 $t=1800~\mathrm s$ 时，
\[
\ell_T\sim\sqrt{\alpha t}
\approx1.74~\mathrm{cm},
\]
已接近药材半径 $R=2~\mathrm{cm}$；
而
\[
\ell_C\sim\sqrt{D(C_0)t}
\approx0.30~\mathrm{cm},
\]
仅约为半径的15\%。
因此有
\[
\ell_T\gg\ell_C,
\]
与数值结果中“中心温度明显升高，而中心含水率基本不变”的现象一致。
预热阶段由此表现出明显的“热量先深入、水分后迁移”特征。

\IfFileExists{问题1_温度时空云图.png}{
\begin{figure}[H]
\centering
\includegraphics[width=0.82\textwidth]{问题1_温度时空云图.png}
\caption{问题一温度场时空分布}
\end{figure}
}{}

\IfFileExists{问题1_水分浓度时空云图.png}{
\begin{figure}[H]
\centering
\includegraphics[width=0.82\textwidth]{问题1_水分浓度时空云图.png}
\caption{问题一水分浓度场时空分布}
\end{figure}
}{}

\subsection{模型讨论与问题二衔接}

问题一的主要局限在于：
一维模型忽略端面附近的二维效应；
附件1中的水分浓度采用等效 Robin 边界解释；
同时暂未考虑蒸发潜热反馈。
此外，本问固定圆柱几何尺寸，失水收缩将在问题四中进一步处理。

问题一由此建立了固定几何、常热物性的基础模型。
完整干燥过程中，含水率持续下降会使
$\rho$、$c_p$、$k$ 及 $D$ 发生变化，
因此问题二根据附录3解除常物性假设，
在同一径向有限体积框架上建立变物性热质耦合模型：
\[
\boxed{
M_1:\ \text{常物性基础模型}
}
\quad\Longrightarrow\quad
\boxed{
M_2:\ \text{变物性热质耦合模型}
}.
\]

\begin{thebibliography}{9}

\bibitem{wang2024}
王乐意, 李长河, 刘明政, 等.
中药材干燥技术与装备研究现状[J].
农业工程学报, 2024, 40(2): 1--28.
DOI: 10.11975/j.issn.1002-6819.202306104.

\bibitem{dasilva2014}
DA SILVA W P, E SILVA C M D P S, GAMA F J A.
Estimation of thermo-physical properties of products with cylindrical shape during drying:
The coupling between mass and heat[J].
Journal of Food Engineering, 2014, 141: 65--73.
DOI: 10.1016/j.jfoodeng.2014.05.010.

\end{thebibliography}

\end{document}
